\section{Discussion and Future Work}
\label{sec:discussion}

% Target: ~0.75 two-column page (~1.5 single-column equivalent)
% Condense report sections/future_work.tex.

\subsection{Hardware Verification}
\label{sub:hw_verification}

% [~0.3 column]
%   - QEMU evaluation demonstrates constant-time execution paths but does
%     not validate microarchitectural isolation. fence.t behaviour (cache
%     line flush, TLB invalidation, branch-predictor reset) is not modelled.
%   - Next step: port to a cycle-accurate simulator (gem5 with RISC-V
%     support) or real silicon with fence.t. A gem5 evaluation would allow
%     measuring actual microarchitectural state leakage and validating that
%     fence.t closes all on-core channels.
%   - Preliminary attempt at a QEMU TCG plugin for cycle-accurate modelling
%     was unsuccessful due to the complexity of emulating cache and pipeline
%     behaviour at the instruction level.

\subsection{System-Call Audit}
\label{sub:syscall_audit}

% [~0.3 column]
%   - Not all FreeRTOS system calls are guaranteed constant-time with
%     respect to domain state. Calls that acquire the scheduler lock
%     (`vTaskSuspendAll`) create critical sections that delay domain
%     switches.
%   - Proposed mitigation: a two-level locking system distinguishing
%     cross-domain critical sections (block domain switches) from
%     within-domain critical sections (only block intra-domain preemption).
%   - Example: the S3K TTAS (test-and-test-and-set) lock design could
%     serve as a model for constant-time cross-domain synchronisation.
%   - A full audit of all FreeRTOS system calls remains future work.

\subsection{Remaining Timing Channels}
\label{sub:remaining_channels}

% [~0.5 column]
%   - External interrupts: not partitioned by domain. An interrupt handler
%     executes in the currently active domain's context, creating a timing
%     channel.
%   - Higher-level caches (L2, L3): not flushed by fence.t. Cache colouring
%     (as in seL4) or hardware cache partitioning (Intel CAT, ARM MPAM,
%     RISC-V QoS) could address this, but are not implemented.
%   - The S3K approach of placing the kernel in scratchpad memory eliminates
%     kernel cache footprints entirely. This could be adapted to FreeRTOS
%     on hardware with SPM support.
%   - Memory bus contention (DRAM refresh, DMA): beyond the scope of
%     OS-level mitigation.

\subsection{Lighter-Weight Alternatives}
\label{sub:alternatives}

% [~0.3 column]
%   - Full domain scheduling may be overkill for systems that only require
%     isolation between specific ``secure'' and ``untrusted'' tasks.
%   - Alternative: secure-task partitioning. Mark a subset of tasks as
%     security-sensitive. On entry/exit of a secure task, execute fence.t
%     to flush the residual state of the previous (potentially untrusted)
%     task. No time-slice partitioning required.
%   - Alternative: scratchpad-resident buffers. For tasks that compute on
%     secret-indexed data, place buffers in scratchpad memory (as in S3K)
%     to prevent cache footprints from being observable.
%   - These approaches offer a spectrum of protection vs.\ overhead; the
%     domain scheduler represents the strongest but most costly point on
%     this spectrum.
